\documentclass[10pt]{article}

\usepackage[utf8]{inputenc}

\usepackage[T1]{fontenc}

\usepackage{amsmath}

\usepackage{amsfonts}

\usepackage{amssymb}

\usepackage[version=4]{mhchem}

\usepackage{stmaryrd}

\usepackage{graphicx}

\usepackage[export]{adjustbox}

\graphicspath{ {./images/} }



\begin{document}

каждой прямой существует такая скрещивающаяся с ней прямая, что полюсы плоскостей, проходящих через одну прямую, лежат на другой, а полярные плоскости точек, лежащих на одной прямой, проходят через другую. Такие скрещивающиеся прямые наз. в 3 а и м н о п о л я р ны м и (взаимные поляры). Две скрещивающиеся прямые наз. к о с ы м и, если они не являются взаимными полярами или если каждая из них не пересекает поляру другой. Две косые прямые имегт два общих перпендикуляра, являющихся взаимными полярами. Если два общих перпендикуляра косых прямых неравновелики (имеют разные длины), то әти косые прямые наз. р а с х о д я щ и м и с я, а длины общих перпендикуляров дают наименьшее и наибольшее расстояния одной прямой от другой. Две косые прямые, имеющие бесконечное множество общих перпендикуляров одинаковой длины, наз. п а р а л л е л я м и Клиффорда (ра в ноотстоя щи и и, пара т а к т ич н м и пр я мым и). Через каждую точку пространства, лежащую вне данной прямой и вне полярной ей прямой, проходят две параллели Клиффорда к данной прямой. Множество точек, отстоящих от данной прямой на одном и том же расстоянии, меньшем $\pi r / 2$, наз. по в е $\mathbf{p} \mathbf{x}$ о с т ь ю Данная прямая наз. осью, а расстояние $\rho$ точек от оси радиусом поверхности Клиффорда. Эта поверхность имеет две взаимно полярные оси и соответственно два радиуса, дополняющих друг друга до полупрямой. Плоскости, проходящие через ось поверхности Клиффорда, пересекают ее по окружности. Через каждую точку поверхности Клиффорда проходят две прямые, равноотстоящие от ее осей и целиком принадлежащие поверхности; эти прямые наз. п р я м о ли н е й ными образ ующимипов прхвости ф о р д а. Всякие три прямые, равноотстоящие друг от друга, определяют поверхность Клиффорда, для к-рой они являются образующими. Каждая пара прямолинейных образующих различных семейств пересекается под постоянным углом. Поверхность Клиффорда изометрична евклидову ромбу с острым углом, равным углу между образующими различных семейств, и стороной длины $\pi r$, у к-рого отождествлены точки противоположных сторон, соединяемых прямыми, параллельными другим сторонам. Иными словами, поверхность Клиффорда несет на себе евклидову геометрию. Площадь поверхности Клиффорда радиуса $\rho$ равна $\pi^{2} r^{2} \sin (2 \rho / r)$.



С ф е р а в $\partial^{3}$ - множество точек, расположенных на одинаковом расстоянии от нек-рой точки (ее ц е п тр a). Сфера является эквидистантой нек-рой плоскости (б а зы с ф е р ы). Площадь сферы радиуса $R$ равна $4 \pi r^{2} \sin ^{2}(R / r)$.



В Р. г. пространства имеет место п р и н ц и п д в о йс т в е н н о с т и: во всяком верном утверждении можно взаимно заменить термины «точка» и «плоскость», в результате получается также верное утверждение.



По-видимому, первое сообщение о Р. г. сделано Б. Риманом (B. Riemann) в его лекции "О гипотезах, лежащих в основании геометрии» (1854, опубл. 1867, см. [1]), где Р. г. рассматривалась как риманова геометрия постоянной положительной гауссовой кривизны.



Лum.: [1] Р и м а н Б., в кн.: Об основаниях геометрии, М., 1956, с. 309-25; [2] Е фи м о в Н. В., Высшая геометрия, 6 изд., М., 1978; [3] Р о 3 е и ф е л ь д Б. А., Неевклидовы прост-



\includegraphics[max width=\textwidth, center]{2023_07_08_7e037e27edba11ae4eeeg-1}

ч. 2, М., 1956; 15] Б о г о м о Л. в С. А., Введение в неевклидо-



РИМАНА ГИПОТЕЗЫ в $а$ н а л и т и ч е с к о й т е о р и и ч и с е л-пять гипотез, высказаныых Б. Риманом (B. Riemann, 1876) относительно распределения нетривиальных нулей дзета-функции $\zeta(s)$ и относительно выражения через эти нули числа простых чисел, не превосходящих $x$ (см. Дзета-бункция). Не доказана и не опровергнута одна Р. г.: все нетривиальные нули дзета-функции $\zeta(s)$ лежат на прямой $\operatorname{Re} s=1 / 2$. A. Ф. Лаврur. РИМАНА ДЗЕТА-ФУНКЦИЯ - см. Дзета-функция. РИМАНА ДИФФЕРЕНЦИАЛЬНОЕ УРАВНЕНИЕ линейное однородное обыкновенное дифффе ренциальное уравнение 2-го порядка в комплексной плоскости, имеющее три заданные регулярные особые точки $a, b, c$ с соответствующими характеристич. показателями $\alpha$, $\alpha^{\prime}, \beta, \beta^{\prime}, \gamma, \gamma^{\prime}$ в этих точках. Общий вид такого уравнения впервые выписал Э. Паппериц (Е. Рapperitz), из-за чего оно также наз. Папперица уравнением. Решения P. д. у. записываются в виде так наз. $P$-ф у н к ц и и Р и м а на



\[

w=P\left\{\begin{array}{lll}

a & b & c \\

\alpha & \beta & \gamma \\

\alpha^{\prime} & \beta^{\prime} & \gamma^{\prime}

\end{array}\right\} .

\]



Р. Д. У. принадлежит Фукса классу уравнений с тремя особыми точками. Частным случаем Р. д. У. является гипергеометрическое уравнение (особые точки: $0,1, \infty)$; поэтому само P. Д. У. иногда наз. о б о б щ е н ны м гиперге пмет рическия уравнением. Р. д. у. приводится к Похгамлера уравнению, а потому решение Р. д. у. можно записать в виде интеграла по специальному контуру в комплексной плоскости.



Лит. см. при ст. Папперица уравнение. Н. Х. Розов.



РИМАНА ИНТЕГРАЈ - обобщение понятия Коши интеграла на нек-рый класс разрывных функций, введенное Б. Риманом (B. Riemann, 1853). Пусть функция $f(x)$ задана на отрезке $[a, b]$ и $a=x_{0}<x_{1}<\ldots<x_{i-1}<$ $<x_{i}<\ldots<x_{n}=b, \quad i=1,2, \ldots, n, x_{i}-x_{i-1}=\Delta x_{i}$. Сумму вида



\[

\begin{gathered}

\sigma=f\left(\xi_{1}\right)\left(x_{1}-x_{0}\right)+\ldots+f\left(\xi_{i}\right)\left(x_{i}-x_{i-1}\right)+\ldots \\

\ldots+f\left(\xi_{n}\right)\left(x_{n}-x_{n-1}\right),

\end{gathered}

\]



где $x_{i-1} \leqslant \xi_{i} \leqslant x_{i}$, наз. и н т е г р а ль н о й с у м мой, отвечающей данному разбиению отрезка $[a, b]$ точками $x_{i}$ и выбору точек $\xi_{i}$. Число $I$ наз. пределом интегральных сумм (1) при $\max \Delta x_{i}<\delta$, если для любого $\varepsilon>0$ найдется $\delta>0$ такое, что при $\underset{i}{ } \max _{i} \Delta x_{i} \rightarrow 0$ справедливо неравенство $|\sigma-I|<\varepsilon$. Если существует конечный предел $I$ интегральных сумм при $\underset{i}{\max } \Delta x_{i} \rightarrow 0$, то функцию $f(x)$ наз. ин т е г $\mathbf{p}$ и р у е м ой в с мыс ле $\mathrm{P}$ иман а на отрезке $[a, b]$ при $a<b$, а указанный предел олред еленным интегра лом Р имава от функции $f(x)$ по отрезку $[a, b]$ и обозначают



\[

\int_{a}^{b} f(x) d x

\]



При $a=b$, по определению, полагают



\[

\int_{a}^{a} f(x) d x=0

\]



а при $a>b$ определяют интеграл (2) с помощью равенства



\[

\int_{a}^{b} f(x) d x=-\int_{b}^{a} f(x) d x

\]



Необходимым и достаточным условием интегрируемости $f(x)$ на $[a, b]$ в смысле Римана являются ограниченность $f(x)$ на этом отрезке и равенство нулю Лебега меры множества всех точек разрыва $f(x)$, содержащихся на $[a, b]$.



С в о й с т в а Р. и. 1) Всякая интегрируемая по Риману на отрезке $[a, b]$ функция $f(x)$ ограничена на этом отрезке (обратное неверно: примером ограниченной и неинтегрируемой на $[a, b]$ функции служит Дирихле бункция).





\end{document}